\documentclass[11pt]{article}
\usepackage{fontspec}
\usepackage[french,english]{babel}
\usepackage{amsmath}
\usepackage{amssymb}
\usepackage{amsthm}
\usepackage{mathrsfs}
\usepackage{tikz-cd}
\usepackage{enumitem}

\newcommand{\printedpage}[1]{\marginpar{\footnotesize p.~#1}}
\newcommand{\msp}[1]{\marginpar{\footnotesize [#1]}}
\newcommand{\sDia}{\mathsf{Dia}}
\newcommand{\sOb}{\mathsf{Ob}}
\newcommand{\sid}{\mathsf{id}}
\newcommand{\DD}{\mathbb{D}}

\begin{document}

\printedpage{19}
% head: A. Grothendieck : Dérivateurs, Ch. I, Édition des Universités de Montpellier II et Paris VII
% foot: Version 8 juillet 2001
\msp{page 28}
that the functor for the associated fundamental categories
\[ \mathcal{M}(I) \xrightarrow{f^{\mathcal{M}}_*} \mathcal{M}(I') \]
is right exact. Now here it is a direct image functor of ordinary presheaves, and this functor has no desire at all to be \textit{right} exact. (Even if $\mathcal{M}$ is \textit{abelian}, for that matter!) Let us take for example for $f : I \hookrightarrow I'$ an open immersion with $I' = \boldsymbol{\Delta}_1 \times \boldsymbol{\Delta}_1$, $I = I' - \{c\}$, $c$ a final object of $I'$
\[
I' = \begin{tikzcd}
c & b \arrow[l] \\
a \arrow[u] & 0 \arrow[l] \arrow[u]
\end{tikzcd}
\qquad\qquad
I = \begin{tikzcd}
 & b \\
a & 0 \arrow[l] \arrow[u]
\end{tikzcd}
\]
The functor $f_*$ is the one which associates to a diagram
\[
F = \left\{
\begin{tikzcd}
 & F_b \arrow[d] \\
F_a \arrow[r] & F_0
\end{tikzcd}
\right\}
\]
the completed cartesian square
\[
\overline{F} = \left\{
\begin{tikzcd}
F_a \times_{F_0} F_b \arrow[r] \arrow[d] & F_b \arrow[d, "\beta"] \\
F_a \arrow[r, "\alpha"] & F_0
\end{tikzcd}
\right\}
\]
Now if $F \mapsto \overline{F}$ were right exact, $F \mapsto F_a \times_{F_0} F_b$ would be so too. Now $F_a \times_{F_0} F_b = K(F) = \mathsf{Ker}(F_a \times F_b \xrightarrow{\alpha pr_1 - \beta pr_2} F_0)$, as an additive functor in $F$, is not right exact, for example it does not necessarily transform epimorphism into epimorphism. Let us take for example
\[
F' = \left\{
\begin{tikzcd}
 & F_b \arrow[d] \\
F_a \arrow[r] & 0
\end{tikzcd}
\right\}
\]
and the obvious homomorphism $F \to F'$ which is the identity on
\msp{page 29}
$F_a$ and on $F_b$. It is an epimorphism and the corresponding homomorphism $\overline{F} \to \overline{F'}$ is the inclusion
\[ F_a \times_{F_0} F_b \longrightarrow F_a \times F_b \]
which is an epimorphism only if $\alpha = \beta = 0$.

\printedpage{20}
% head: A. Grothendieck : Dérivateurs, Ch. I, Édition des Universités de Montpellier II et Paris VII
% foot: Version 8 juillet 2001
\textbf{Perplexity} : One should be able to define a notion of an (additive) ``\textit{triangulated}'' derivator (which would correspond to the intuition of triangulated categories; it is a derivator satisfying fairly strong axioms), so that for morphisms between these, left exactness and right exactness are \textit{equivalent.}

\textbf{``Base change theorem''} for derivators. Let $\DD$ be a left derivator (existence of $f_*$) and let us consider
\[ f : I \longrightarrow J \qquad \text{in } \sDia \quad , \]
let $j_0 \in \sOb\ J$, we propose to ``calculate'' the functor
\[ F \mapsto f_*(F)_{j_0} = \alpha^*_{j_0}(f_*(F)) \quad , \]
where
\[ \alpha_{j_0} : e \longrightarrow J \]
is the inclusion of the one-point category $e$ in $J$, with value $j_0$.

\msp{page 30}
Let us consider the diagram
\[
\begin{tikzcd}[row sep=large, column sep=large]
I/j_0 \arrow[r, "q"] \arrow[d, "f'"'] & I \arrow[d, "f"] \\
e \arrow[r, "p = \alpha_{j_0}"'] & J
\arrow[from=1-2, to=2-1, dashed, shorten <=14pt, shorten >=14pt, "\gamma"']
\end{tikzcd}
\quad ,
\tag*{(1)}
\]
where $I/j_0$ is the category $I \times_J J/j_0$ of pairs
\[ \{(i,\varphi) \mid i \in \sOb\ I,\ \varphi : f(i) \to j_0\} \]
and $q$ is the functor $(i,\varphi) \mapsto i$. This diagram is not commutative, but there is a canonical functor \texttt{[rather morphism of functors]}
\[ fq \xrightarrow{\gamma} pf' \quad , \]
where $fq$ is the functor $(i,\varphi) \mapsto f(i)$ and $pf'$ the constant functor $(i,\varphi) \mapsto j_0$ and where $\gamma(i,\varphi) : f(i) \to j_0$ is defined by $\gamma(i,\varphi) = \varphi$.

More generally, let us consider a diagram in $\sDia$
\[
\begin{tikzcd}[row sep=large, column sep=large]
I' \arrow[r, "q"] \arrow[d, "f'"'] & I \arrow[d, "f"] \\
J' \arrow[r, "p"'] & J
\arrow[from=1-2, to=2-1, dashed, shorten <=14pt, shorten >=14pt, "\gamma"']
\end{tikzcd}
\quad ,
\tag*{(2)}
\]
\printedpage{21}
% head: A. Grothendieck : Dérivateurs, Ch. I, Édition des Universités de Montpellier II et Paris VII
% foot: Version 8 juillet 2001
with $\gamma : fq \to pf'$. We then define a \textit{base change morphism for direct images} relative to this diagram
\[ p^* f_* \xrightarrow{c_\gamma} f'_* q^* \]
\msp{page 31}
in two ways, which give the same result (general story of adjoint functors to be made explicit once and for all, in terms of a diagram of functors
\[
\begin{tikzcd}[row sep=large, column sep=large]
X' & X \arrow[l, "q^*"'] \\
Y' \arrow[u, "{f'}^*"] & Y \arrow[l, "p^*"'] \arrow[u, "f^*"']
\arrow[from=2-1, to=1-2, dashed, shorten <=14pt, shorten >=6pt]
\end{tikzcd}
\]
assuming that the left adjoint functors $f'_*$, $f_*$ of the two vertical functors exist). An interesting case is the one where
\[ pf' = fq \quad , \qquad \gamma = \sid_{fq} \]
and where we put $c_\gamma = c$.

Here
\[ p^* f_* \xrightarrow{?} f'_* q^* \]
amounts to
\[
\begin{tikzcd}[column sep=small]
f_* \arrow[r, "?"] & p_* f'_* q^* \simeq & (pf')_* q^* \\
 & & (fq)_* q^* \simeq (f_* q_*) q^* = f_*(q_* q^*) \quad . \arrow[u, "\gamma_* * q^*"']
\end{tikzcd}
\]
Now we have
\[ \sid \xrightarrow{\alpha} q_* q^* \qquad\qquad \text{(adjunction)} \quad , \]
whence
\[ f_* \xrightarrow{f_* * \alpha} f_*(q_* q^*) \]
and composing with $\gamma * q^*$, we find what we want. And instead of this we can look for
\[
\begin{tikzcd}
{f'}^*(p^* f_*) \arrow[r, "?"] \arrow[d, equal] & q^* \\
(pf')^* f_* \arrow[d, "\gamma^* * f_*"] & \\
(fq)^* f_* \simeq q^*(f^* f_*) \arrow[r, "q^* * \gamma'"] & q^*
\end{tikzcd}
\quad ,
\]
where $\gamma' : f^* f_* \to \sid$ is the adjunction.

\printedpage{22}
% head: A. Grothendieck : Dérivateurs, Ch. I, Édition des Universités de Montpellier II et Paris VII
% foot: Version 8 juillet 2001
\msp{page 32}
\textbf{Der}$_4$ : The base change morphism for direct images $f_*$ associated with a diagram (1) is always an isomorphism (left derivators). Dual condition for right derivators and base change for direct images $f_!$, for the diagram
\[
\begin{tikzcd}[row sep=large, column sep=large]
j_0 \backslash I \arrow[r] \arrow[d] & I \arrow[d, "f"] \\
e \arrow[r, "p = \alpha_{j_0}"'] & J
\arrow[from=2-1, to=1-2, dashed, shorten <=22pt, shorten >=8pt, "\gamma"]
\end{tikzcd}
\tag*{(1$'$)}
\]
associated with $f : I \to J$ in $\sDia$, and $j_0 \in \sOb\ J$.

\textbf{Der}$_5$\footnote[10]{\underline{Totally} botched, becomes reasonable only for constant coefficients on $F$.} : Let $f : I \to J$ be a \textit{strict Hot-equivalence, i.e.} $\forall j_0 \in \sOb\ J$, $I/j_0$ is contractible. Then $f^*$ is fully faithful, \textit{i.e.} the adjunction morphism
\[ \sid_{\DD(J)} \longrightarrow f_* f^* \]
is an isomorphism ($\DD$ a left derivator). Dually, if $f$ is a co-strict equivalence, \textit{i.e.} $f^\circ : I^\circ \to J^\circ$ is a strict equivalence, \textit{i.e.} $\forall j_0 \in \sOb\ J$, $j_0 \backslash I$ is contractible, then (if $\DD$ is a right derivator) the functor $f^*$ is fully faithful, \textit{i.e.} the adjunction morphism % Resolved: the edition line-break gives co-stricte; retain that spelling.
\[ f_! f^* \longrightarrow \sid_{\DD(J)} \]
is an isomorphism.

N.B. Thus the functors $f_*$ (resp. $f_!$) are localization functors.

\textbf{N.B.} Pages 33 to 40 have been deleted, because botched. % [sic] suprimées -> supprimées, spelling

\msp{page 41}
I am going to take up again the terminology and the axioms relating to derivators.

A \textit{prederivator} is a 2-functor
\[ \sDia^\circ \xrightarrow{\DD} \mathcal{C}at \quad , \]
where $\sDia$ is a full 2-subcategory of $\mathcal{C}at$.\footnote[11]{$\sDia$ \textit{domain} of the derivator. The objects are the \textit{index categories} for $\DD$. The $\DD(I)$ are the categories of coefficients for $\DD$, the objects of $\DD(I)$ the ``coefficients of type $\DD$ on $I$''. % [sic] coeficients -> coefficients, spelling
} We assume at least that $\sDia$ contains the \textit{finite ordered sets,} that $\sDia$ is stable under finite $\varprojlim$, under finite sums. That with every category $C$ and every strictly full open (resp. closed) subcategory it contain the complementary subcategory.\footnote[12]{Notation $f^*$ : inverse image of coefficients.}

\printedpage{23}
% head: A. Grothendieck : Dérivateurs, Ch. I, Édition des Universités de Montpellier II et Paris VII
% foot: Version 8 juillet 2001
\msp{page 42}
We define the prederivator $\DD_{\mathcal{M},\Sigma}$ associated with a system
\[ (\mathcal{M},\Sigma) \qquad\qquad \mathcal{M} \text{ Category,} \quad \Sigma \subset \mathsf{Fl}(\mathcal{M}) \quad . \] % Resolved: Fl abbreviates fleches, with lowercase l; retain Fl.
If $\Sigma = \emptyset$, we write $\DD_{\mathcal{M}}$.

Prederivators form a 2- category $\mathsf{PREDER}$, and
\[ (\mathcal{M},\Sigma) \longmapsto \DD_{\mathcal{M},\Sigma} \]
is a 2-functor
\[ \mathsf{PREMOD}
\begin{pmatrix}
\text{``Category of categories} \\
\text{of models'' (weak)}
\end{pmatrix}
\longrightarrow \mathsf{PREDER} \quad . \]
For $\DD$ given, we define the induced derivator ($I \in \sDia$)
\[ \DD_I = \big(J \longmapsto \DD(I \times J)\big) \quad , \]
depending on $I$ contravariantly, whence a 2-functor
\[ \begin{aligned}
\sDia^\circ &\longrightarrow \mathsf{PREDER} \\
I &\longmapsto \DD_I
\end{aligned} \]
or again a bi-2-functor
\[ \begin{aligned}
\mathsf{PREDER} \times \sDia^\circ &\longrightarrow \mathsf{PREDER} \\
(\DD, I) &\longmapsto \DD_I \quad .
\end{aligned} \]
If $\DD$ is defined by $(\mathcal{M},\Sigma)$, then $\DD_I$ is defined by $(\mathcal{M}(I),\Sigma(I))$ :
\[ (\DD_{\mathcal{M},\Sigma})_I \simeq \DD_{\mathcal{M}(I),\Sigma(I)} \quad , \]
\msp{page 43}
where $\mathcal{M}(I) = \underline{\mathsf{Hom}}(I^\circ,\mathcal{M})$ (presheaves on $I$ with values in $\mathcal{M}$), and $\Sigma(I) \subset \mathsf{Fl}\ \mathcal{M}(I)$ is deduced from $\Sigma$ by looking fiber by fiber (now assuming $\Sigma$ saturated). % Resolved: the repeated arrow-class symbol is Fl, with lowercase l.

For any axioms we shall single out for certain prederivators, we shall have to check that they are stable under passage to the induced prederivators $\DD_I$, and also how they behave under duality (passage from $\DD$ to $\DD^\circ$). % [sic] qu'il sont -> qu'ils sont, agreement

Definition of the dual prederivator $\DD^\circ$
\[ \DD^\circ(I) \overset{\text{def}}{=} \DD(I^\circ)^\circ \qquad\qquad \textit{etc.} \]
\textbf{Der 1} (First localization axiom) : \textit{a}) If $I$, $J \in \sDia$, then
\[ \DD(I \amalg J) \longrightarrow \DD(I) \times \DD(J) \]
is an equivalence of categories.

\textit{b}) $\DD(\emptyset)$ is the final category.

Stable under passage to the induced prederivators, and passage to the dual.

\textbf{Corollary.} \textit{Let} $(I_\alpha)_{\alpha \in A}$ be \texttt{[a]} \textit{finite family of categories in} $\sDia$, $I = \coprod\limits_{\alpha \in A} I_\alpha$. \textit{Then}
\[ \DD(I) \longrightarrow \prod_{\alpha \in A} \DD(I_\alpha) \]
\textit{is an equivalence of categories.}

\end{document}
